% Bagian: second-order-logic
% Bab: syntax-and-semantics
% Subbagian: size-of-domain

\documentclass[../../../include/open-logic-section]{subfiles}

\begin{document}

\olfileid[id]{sol}{syn}{siz}

\olsection{Mendeskripsikan \usetoken{P}{domain} Tak Hingga dan
  \usetoken{S}{enumerable}}

Suatu himpunan~$M$ bersifat tak hingga Dedekind jika dan hanya jika
terdapat fungsi !!a{injective}~$f\colon M \to M$ yang tidak
!!{surjective}, yakni, dengan $\ran{f} \neq M$. Dalam logika orde pertama,
kita dapat mempertimbangkan suatu !!{function} beraritas satu~$f$ dan
menyatakan bahwa fungsi~$\Assign{f}{M}$ yang ditetapkan kepadanya dalam
!!a{structure}~$\Struct{M}$ bersifat !!{injective} dan
$\ran{f} \neq \Domain{M}$:
\[
\lforall[x][\lforall[y][(\eq[f(x)][f(y)] \to \eq[x][y])]] \land
\lexists[y][\lforall[x][\eq/[y][f(x)]]].
\]
Jika $\Struct{M}$ memenuhi !!{sentence} ini, maka
$\Assign{f}{M}: \Domain{M} \to \Domain{M}$ bersifat !!{injective},
sehingga $\Domain{M}$ harus tak hingga. Jika $\Domain{M}$ tak hingga dan
karena itu terdapat fungsi semacam itu, kita dapat mengambil
$\Assign{f}{M}$ sebagai fungsi tersebut, dan $\Struct{M}$ akan memenuhi
!!{sentence} itu. Namun, cara ini mengharuskan bahasa kita memuat simbol
nonlogis~$f$ yang kita gunakan untuk tujuan tersebut. Dalam logika orde kedua,
kita cukup menyatakan bahwa fungsi semacam itu \emph{ada}. Hal ini tidak lagi
memerlukan~$f$, dan kita memperoleh !!{sentence} dalam logika orde kedua murni
\[
\fn{Inf} \ident \lexists[u][(\lforall[x][\lforall[y][(\eq[u(x)][u(y)]
      \lif \eq[x][y])]] \land \lexists[y][\lforall[x][\eq/[y][u(x)]]])].
\]
$\Sat{M}{\fn{Inf}}$ jika dan hanya jika $\Domain{M}$ tak hingga. Selanjutnya,
kita dapat mendefinisikan $\fn{Fin} \ident \lnot \fn{Inf}$;
$\Sat{M}{\fn{Fin}}$ jika dan hanya jika $\Domain{M}$ berhingga. Tidak ada satu
!!{sentence} logika orde pertama murni yang dapat mengungkapkan bahwa
!!{domain} tak hingga, meskipun suatu himpunan tak hingga kalimat semacam itu
dapat mengungkapkannya. Tidak ada himpunan !!{sentence}s logika orde pertama murni
yang terpenuhi dalam !!a{structure} jika dan hanya jika domainnya berhingga.

\begin{prop}
$\Sat{M}{\fn{Inf}}$ jika dan hanya jika $\Domain{M}$ tak hingga.
\end{prop}

\begin{proof}
$\Sat{M}{\fn{Inf}}$ jika dan hanya jika
  $\Sat{M}{\lforall[x][\lforall[y][(\eq[u(x)][u(y)] \lif \eq[x][y])]]
    \land \lexists[y][\lforall[x][\eq/[y][u(x)]]]}[s]$ bagi suatu~$s$.
Jika hal itu berlaku, $s(u)$ merupakan fungsi !!a{injective}, dan suatu
$y \in \Domain{M}$ tidak berada dalam daerah hasil~$s(u)$. Sebaliknya, jika
terdapat !!a{injective} $f\colon \Domain{M} \to \Domain{M}$ dengan
$\ran{f} \neq \Domain{M}$, maka suatu penugasan variabel dengan
$s(u) = f$ menjadi saksi yang diperlukan.
\end{proof}

Suatu himpunan tak kosong $M$ bersifat !!{enumerable} jika terdapat suatu
enumerasi
\[
m_0, m_1, m_2, \dots
\]
atas !!{element}s himpunan tersebut (tanpa pengulangan, tetapi mungkin
berhingga). Enumerasi
semacam itu ada jika dan hanya jika terdapat !!a{element}~$z \in M$ dan suatu
fungsi $f\colon M \to M$ sedemikian sehingga $z$, $f(z)$, $f(f(z))$, \dots,
mencakup semua !!{element}s dari~$M$. Sebab, jika enumerasi itu ada, ambil
$z = m_0$ dan $f(m_k) = m_{k+1}$ (atau $f(m_k) = m_k$ jika $m_k$ merupakan
!!{element} terakhir dari enumerasi); itulah !!{element} dan fungsi yang
diperlukan. Sebaliknya, jika $z$ dan $f$ semacam itu ada, maka dengan
menghapus pengulangan dari $z$, $f(z)$, $f(f(z))$, \dots{} diperoleh suatu
enumerasi atas~$M$, sehingga $M$ bersifat !!{enumerable}. Kita dapat
mengungkapkan keberadaan $z$ dan~$f$ dalam logika orde kedua untuk menghasilkan
suatu !!{sentence} yang benar dalam !!a{structure} jika dan hanya jika domain
!!{structure} tersebut !!{enumerable}:
\[
\fn{Count} \ident
\lexists[z][\lexists[u][\lforall[X][((X(z) \land
      \lforall[x][(X(x) \lif X(u(x)))]) \lif \lforall[x][X(x)])]]]
\]

\begin{prop}
$\Sat{M}{\fn{Count}}$ jika dan hanya jika $\Domain{M}$ !!{enumerable}.
\end{prop}

\begin{proof}
Andaikan $\Domain{M}$ !!{enumerable}, dan misalkan $m_0$, $m_1$, \dots,
merupakan suatu enumerasi. Dengan menghapus pengulangan, kita dapat menjamin
bahwa tidak ada $m_k$ yang muncul dua kali. Definisikan
$f(m_k) = m_{k+1}$ bagi setiap anggota yang bukan anggota terakhir; jika
enumerasinya berhingga, petakan anggota terakhir itu ke dirinya sendiri.
Lalu, tetapkan $s(z) = m_0$ dan $s(u) = f$. Kita tunjukkan bahwa
\[
\Sat{M}{\lforall[X][((X(z) \land \lforall[x][(X(x) \lif X(u(x)))])
    \lif \lforall[x][X(x)])]}[s]
\]
Misalkan $N \subseteq \Domain{M}$ sebarang. Andaikan pula bahwa
$\Sat{M}{(X(z) \land \lforall[x][(X(x) \lif
X(u(x)))])}[\Subst{s}{N}{X}]$. Maka $\Subst{s}{N}{X}(z) \in N$, dan setiap
kali $a \in N$, berlaku pula $(\Subst{s}{N}{X}(u))(a) \in N$. Dengan kata
lain, karena $\varAssign{\Subst{s}{N}{X}}{s}{X}$, $m_0 \in N$ dan jika
$a \in N$, maka $f(a) \in N$; jadi $m_0 \in N$,
$m_1 = f(m_0) \in N$, $m_2 = f(f(m_0)) \in N$, dan seterusnya. Dengan
demikian, $N = \Domain{M}$, sehingga
$\Sat{M}{\lforall[x][X(x)]}[\Subst{s}{N}{X}]$. Karena
$N \subseteq \Domain{M}$ dipilih secara sebarang, kita selesai:
$\Sat{M}{\fn{Count}}$.

Sekarang andaikan $\Sat{M}{\fn{Count}}$, yakni,
\[
\Sat{M}{\lforall[X][((X(z) \land \lforall[x][(X(x) \lif X(u(x)))])
    \lif \lforall[x][X(x)])]}[s]
\]
bagi suatu~$s$. Misalkan $m = s(z)$ dan $f = s(u)$, lalu perhatikan
$N = \{m, f(m), f(f(m)), \dots\}$. Himpunan $N$ yang didefinisikan demikian
jelas !!{enumerable}. Berdasarkan asumsi,
\[
\Sat{M}{(X(z) \land \lforall[x][(X(x) \lif X(u(x)))])
    \lif \lforall[x][X(x)]}[\Subst{s}{N}{X}]
\]
Selain itu, $\Sat{M}{X(z)}[\Subst{s}{N}{X}]$, karena
$N \ni m = \Subst{s}{N}{X}(z)$, dan juga
$\Sat{M}{\lforall[x][(X(x) \lif X(u(x)))]}[\Subst{s}{N}{X}]$, karena setiap
kali $a \in N$, berlaku pula $f(a) \in N$. Jadi, karena anteseden maupun
kondisional tersebut terpenuhi, konsekuennya juga harus terpenuhi:
$\Sat{M}{\lforall[x][X(x)]}[\Subst{s}{N}{X}]$. Namun, ini berarti bahwa
$N = \Domain{M}$, sehingga $\Domain{M}$ !!{enumerable}, sebab berdasarkan
definisi $N$ bersifat demikian.
\end{proof}

\begin{prob}
!!{sentence} $\fn{Inf} \land \fn{Count}$ benar tepat pada domain yang
!!{denumerable}. Sesuaikan definisi $\fn{Count}$ agar menjadi
suatu !!{sentence} lain yang secara langsung mengungkapkan bahwa domain
!!{denumerable}, dan buktikan bahwa kalimat itu memang demikian.
\end{prob}

\end{document}
